% Apparatus schematic for §3.3.
%
% Routing rules:
%   - Main flow: straight horizontal/vertical only.
%   - Tool response: ONE right angle. UP from MCP, then LEFT into primary.east.
%     Use TikZ |- operator so the right angle is geometrically guaranteed
%     (no diagonal segments due to mismatched y-coordinates).
%   - Feedback wrap: leaves the BOTTOM of verify, wraps around the LEFT
%     of the diagram, arrives at the TOP of primary. All four corners
%     use a uniform clearance \pad so the halo around the diagram is
%     visually balanced.
%
% Label rules: every arrow gets a consistent verb-form gerund
%   (consults / dispatches / flags / tool response).
\begin{figure}[H]
\centering
\begin{tikzpicture}[
  font=\sffamily\small,
  >={Stealth[length=2.4mm]},
  every node/.style={align=center},
  box/.style={rectangle, rounded corners=2pt, draw=accentDark, line width=0.6pt,
              minimum height=11mm, minimum width=34mm, inner sep=4pt,
              fill=paleGray},
  modelbox/.style={box, fill=accentLight, draw=accent, line width=0.8pt},
  decisionbox/.style={diamond, draw=accent, line width=0.8pt, fill=white,
                      inner sep=2pt, aspect=2.4, minimum width=30mm,
                      minimum height=14mm},
  toolbox/.style={box, fill=white, draw=warmGray},
  arrow/.style={->, draw=accentDark, line width=0.6pt},
  feedback/.style={->, draw=condBprimeE, line width=0.7pt, dashed},
  cleanpath/.style={->, draw=condReactive, line width=0.7pt},
  ret/.style={->, draw=accentDark, line width=0.6pt},
]

% ── Grid coordinates ───────────────────────────────────────────────
\def\colL{0}        % left column x
\def\colM{5.6}      % middle column x
\def\colR{11.2}     % right column x
\def\rowT{4.5}      % top row y
\def\rowC{2.25}     % middle row y
\def\rowB{0}        % bottom row y

% Clearances for the feedback wrap-around halo.
% Three independent knobs so you can tune each side without affecting
% the others. Increase a value to push that side further from the
% diagram; decrease to bring it closer.
\def\padBelow{0.6}   % gap below verify before the path turns left
\def\padLeft{1.2}    % gap left of the leftmost column (this is the LEFT channel)
\def\padAbove{0.6}   % gap above primary before the path turns right and down

% ── Nodes ──────────────────────────────────────────────────────────
\node[box]      (case)      at (\colL,\rowT) {Patient vignette\\(natural language)};
\node[modelbox] (extractor) at (\colL,\rowC) {\textbf{Sonnet extractor}\\\footnotesize $M_e$: case $\to$ \texttt{\{field: value\}}};
\node[box]      (facts)     at (\colL,\rowB) {Patient facts\\(reference)};

\node[modelbox]    (primary) at (\colM,\rowT) {\textbf{Qwen3 primary}\\\footnotesize $\langle$tool\_call$\rangle$ generator};
\node[box]         (call)    at (\colM,\rowC) {Tool call payload\\\footnotesize\texttt{input\_data = \{...\}}};
\node[decisionbox] (verify)  at (\colM,\rowB) {\textbf{verify}\\\footnotesize field-by-field};

\node[toolbox] (mcp) at (\colR,\rowB) {MCP calculator\\(deterministic)};

% ── Main flow (all straight) ───────────────────────────────────────
\draw[arrow] (case)      -- (extractor);
\draw[arrow] (extractor) -- (facts);
\draw[arrow] (case)      -- (primary);
\draw[arrow] (primary)   -- (call);
\draw[arrow] (call)      -- (verify);
\draw[arrow] (facts)     -- node[above, font=\footnotesize\itshape,
                                  text=warmGray, yshift=0.5mm]
                            {consults} (verify);
\draw[cleanpath] (verify) -- node[above, font=\footnotesize\bfseries,
                                    text=condReactive] {dispatches}
                              (mcp);

% ── Tool response: ONE right angle via |- operator ─────────────────
% (mcp.north) |- (primary.east) automatically produces:
%   1. Vertical segment UP from mcp.north to primary.east.y
%   2. Right angle at the corner
%   3. Horizontal segment LEFT to primary.east
% No diagonals are possible because |- guarantees axis-aligned paths.
\draw[ret] (mcp.north) |- (primary.east);
% Label: anchored at center on the right vertical lane.
\node[font=\footnotesize\itshape, text=warmGray,
      anchor=center, rotate=90]
  at (\colR+0.30, \rowC) {tool response};

% ── Feedback wrap (left side, balanced halo with uniform \pad) ─────
% Path corners (all axis-aligned):
%   verify.south  → (verify.south.x, verify.south.y - \pad)   [DOWN]
%                 → (colL - boxhalf - \pad, same y)             [LEFT]
%                 → (same x, primary.north.y + \pad)            [UP]
%                 → (primary.north.x, same y)                   [RIGHT]
%                 → primary.north                              [DOWN]
\coordinate (verifyBelow) at ($(verify.south)+(0,-\padBelow)$);
\coordinate (leftLane)    at ($(case.west)+(-\padLeft,0)$);
\coordinate (topLane)     at ($(primary.north)+(0,\padAbove)$);

\draw[feedback]
  (verify.south)
    -- (verifyBelow)                           % DOWN
    -- (leftLane |- verifyBelow)               % LEFT to the left lane
    -- (leftLane |- topLane)                   % UP along the left lane
    -- (primary.north |- topLane)              % RIGHT to above primary
    -- (primary.north);                        % DOWN into primary.north

% Labels positioned along the path:
%   "flags" — short action label at the start of the dashed line
%             (just to the right of where it leaves verify.south).
%   "feedback re-prompt" — descriptive label on the left vertical lane,
%             mirror-image of "tool response" on the right lane.
\node[font=\footnotesize\bfseries, text=condBprimeE,
      anchor=west]
  at ($(verify.south)+(0.15,-0.3)$) {flags};
\node[font=\footnotesize\itshape, text=condBprimeE,
      anchor=center, rotate=90]
  at ($(leftLane)+(-0.30,-2.25)$) {feedback re-prompt};

\end{tikzpicture}
\caption{Apparatus overview. Right: a clean tool call is dispatched to
MCP. Dashed left: a verifier flag re-prompts the primary.}
\label{fig:mechanism}
\end{figure}
